% Chapter 1

\chapter{Introduction} % Main chapter title

\label{Chapter1} % For referencing the chapter elsewhere, use \ref{Chapter1} 

%----------------------------------------------------------------------------------------

While financial institutions generate limited environmental impacts through their direct operations (such as office energy use and IT operations), they are indirectly responsible for substantial greenhouse gas (GHG) emissions by financing high-emitting firms. These are commonly referred to as “financed emissions” and are classified under Scope 3, Category 15 (“Investments”) downstream emissions \citep{ghg_protocol_corporate_2011}.

The relevance of financed emissions begins with the growing recognition that climate change can be a source of financial instability \citep{bolton_green_2020, esrb_esrb_2021}. Climate change may generate physical risks, arising from severe climate-related damages, and transition risks, arising from potentially disorderly mitigation strategies, regulatory changes, carbon pricing, technological change, or shifts in market expectations. At the portfolio level, measuring financed emissions allows financial institutions to assess their exposure to these risks through the companies they finance.

Moving from the portfolio level to the macroeconomic level, the same logic can be extended to the financial system as a whole. If domestic financial institutions are highly exposed to emissions-intensive sectors, financed emissions may become a source of systemic vulnerability. In this sense, the financed emissions of a country’s financial system can provide an indication of broader macro-financial risks. 

In an effort to establish a transparent and harmonized framework for disclosing financed emissions, the Partnership for Carbon Accounting Financials (PCAF) has developed the \textit{Global GHG Accounting and Reporting Standard for the Financial Industry}. While this standard has become widely adopted for attributing emissions at the asset and portfolio levels, its approach suffers from several inherent limitations \citep{pcaf_global_2022}. First, high-quality data remain scarce for certain asset classes, forcing institutions to rely on estimates or proxies. Second, pervasive double counting can occur, whether across different institutions, across asset classes, or when multiple lenders co-finance a single entity. Consequently, bottom-up financed emissions cannot be reliably aggregated to estimate a nation's total financial footprint. 

Beyond its contribution to GHG emissions, the financial system also contributes to broader environmental degradation by financing activities associated with intensive land use, water consumption, and resource extraction. Assessing these wider environmental pressures provides a more comprehensive view of the financial system’s role in climate change and environmental sustainability, thereby supporting the development of more robust sustainability metrics and policy responses. At the same time, the financial system is exposed to climate- and environment-related risks through investment portfolios that finance activities linked to environmental harm.

Building on the macroeconomic framework developed by \textcite{jondeau_environmental_2026}, this thesis quantifies the environmental footprint of two major financial centers: the United States and the Eurozone. The methodology relies on EXIOBASE, an \textit{environmentally extended multi-regional input-output} (EE-MRIO) database that provides GHG emissions, land use, water use, and material extraction at the country-industry level. By integrating these environmental extensions with macro-financial data, the framework attributes environmental pressures to specific groups of financial institutions based on their exposure to real-economy credit recipients. A primary challenge in this attribution is the complexity of financial intermediation chains, where institutions often lend to one another before capital reaches the real economy. Ignoring these intermediation links may lead to double counting or an incomplete measurement of financial exposures. To address this issue, the framework employs \textit{national financial accounts} alongside \textit{from-whom-to-whom} (FWTW) exposure networks. This makes it possible to trace environmental pressures generated by real-economy agents back to their ultimate financiers across the chain of financial intermediation.

The results reveal structural differences between the financial systems of the two regions, which in turn shape their exposure to climate- and nature-related risks. Banks have consistently remained the main source of financing in the Euro area, whereas the US financial system is more market-based, with several groups of financial institutions contributing to market financing. In terms of credit recipients, the US financial system primarily finances its domestic corporate sector. By contrast, Eurozone financial institutions have increasingly allocated funds to foreign counterparties over the past decade.

This shift has contributed to the Eurozone’s financed emissions surpassing those of the United States since 2016 onwards, under the assumption that outbound funds are directed to both developed and emerging markets. In 2022, GHG emissions attributed to the US financial system amounted to 3,761.45 million tonnes of CO\textsubscript{2}e, compared with 4,082.38 million tonnes of CO\textsubscript{2}e attributed to the Eurozone financial centre. This finding challenges the common perception that the Eurozone financial system is relatively green, a perception partly rooted in the European Union’s climate-policy leadership and extensive sustainable-finance regulatory framework, including the European Green Deal, the EU Emissions Trading System, the EU Taxonomy, and sustainable-finance disclosure rules \citep{merler_silvia_how_2025}. Despite this policy and regulatory leadership position, the results show that Eurozone financial institutions remain substantially exposed to emissions generated through their financing activities, particularly through cross-border exposures.

Compared with each region’s territorial and demand-driven emissions, financed emissions are relatively similar in magnitude in the Euro area, whereas in the United States they amount to only around half of these benchmarks. This contrast points to different channels through which emissions-reduction efforts may be prioritized. In the United States, reducing the carbon footprint of the financial system should be complemented by broader decarbonization across domestic industries. In the Euro area, by contrast, the decarbonization of financial portfolios appears especially important, as financed emissions constitute a substantial component of the region’s overall carbon footprint.


Beyond GHG emissions, the analysis also shows that financial-system-based accounting provides a distinct perspective on broader environmental pressures, including land use, water consumption, and resource extraction. These dimensions offer a more comprehensive view of the environmental footprint of financial systems and their exposure to nature-related risks.

The remainder of this thesis is organized as follows. Chapter 2 reviews the relevant literature. Chapters 3 and 4 introduce the conceptual framework and methodology, respectively. Chapter 5 describes the data used in the empirical analysis. Chapter 6 presents the results, and Chapter 7 concludes.